% hip-node-identity-wallet.tex
% Hanzo Improvement Proposal — Node Identity: the Wallet as the Node.
% Number assigned when this lands in HIPs/.
\documentclass[11pt,a4paper]{article}

\usepackage[utf8]{inputenc}
\usepackage[T1]{fontenc}
\usepackage[margin=1in]{geometry}
\usepackage{amsmath,amssymb,amsthm}
\usepackage{booktabs}
\usepackage{longtable}
\usepackage{array}
\usepackage{enumitem}
\usepackage{xcolor}
\usepackage{microtype}
\usepackage[colorlinks=true,linkcolor=black,citecolor=black,urlcolor=blue]{hyperref}

\newtheorem{definition}{Definition}[section]
\newtheorem{claim}[definition]{Claim}
\newtheorem{remark}[definition]{Remark}

% Long identifiers are the norm here, so let the typewriter font break.
\newcommand{\code}[1]{\texttt{\hyphenchar\font=`\-\relax #1}}
\sloppy
\emergencystretch=3em
\newcommand{\Held}{\mathrm{Held}}
\newcommand{\Ran}{\mathrm{Ran}}
\newcommand{\Conf}{\mathrm{Conf}}
\newcommand{\Vrfy}{\mathsf{Vrfy}}
\newcommand{\Sign}{\mathsf{Sign}}
\newcommand{\addr}{\mathsf{addr}}

\setlist[itemize]{itemsep=2pt,topsep=4pt}
\setlist[description]{itemsep=3pt,topsep=4pt}

\title{\textbf{Node Identity: the Wallet as the Node}\\[6pt]
\large A Hanzo Improvement Proposal}

\author{Hanzo AI}
\date{4 September 2026}

\begin{document}
\maketitle

\begin{center}
\begin{tabular}{@{}ll@{}}
\toprule
Number      & assigned on merge into \code{HIPs/} \\
Type        & Standards Track \\
Category    & Core \\
Status      & Draft \\
Requires    & HIP-0024 (Hanzo sovereign L1), HIP-0026 (IAM), HIP-0027 (KMS),\\
            & HIP-0095 (QoS challenge), HIP-0096 (compute rewards),\\
            & HIP-0097 (\code{did:hanzo:}), HIP-1325 (\code{/v1/node}) \\
\bottomrule
\end{tabular}
\end{center}

\begin{abstract}
\noindent
A Hanzo node is a physical or virtual host that boots a Linux microVM, joins one
k3s cluster, and runs the Hanzo cloud. This proposal specifies the identity of
such a node: an account address on the Hanzo sovereign L1, held by the node's
\code{hanzod}, used to sign what the node says about itself, named in the node's
on-chain registration, and paid the node's rewards.

The proposal is written around one distinction, because every claim the system
can make about a node depends on it. A signature establishes that some party
able to evaluate the signing function chose to sign a message. It does not
establish which machine that party was, nor what software that machine was
running. An attestation rooted in a hardware vendor's certificate chain
establishes the second of those, conditionally, and only for platforms that have
the hardware. We name the difference the \emph{possession--platform gap}
(\S\ref{sec:gap}), state what falls on each side of it, and specify the custody,
join, staking, reward, and lifecycle rules that follow from where a given node
falls.

Three properties are load-bearing. First, a node key that can be copied is an
identity that can be forged; custody is therefore specified as a ladder with a
derived assurance class, not as a claim an operator makes. Second, the identity
key (on the box, signing constantly) and the stake owner (off the box, moving
money) are two roles and two keys, because a design in which the compromised key
can also withdraw the bond has no recovery path. Third, forfeiture on-chain is
restricted to self-proving faults --- principally equivocation, which is exactly
what a copied key produces --- so that slashing needs no quorum and no oracle.

Section~\ref{sec:status} and Appendix~\ref{app:fleet} separate what is
implemented from what is proposed. The measurement format, the boot digest, and
the host-signed evidence type exist in \code{hanzoai/node}
(\code{hanzo-libs/hanzo-attest}); the wallet custody backends exist in
\code{hanzoai/cloud} (\code{apps/wallet}) and \code{hanzoai/cli}
(\code{hanzo wallet}); the validator manager exists in \code{luxfi/standard}. The
identity binding between them, the epoch statement, and forfeiture do not exist
and are proposed here.
\end{abstract}

\tableofcontents

% =====================================================================
\section{Scope}

This document specifies the identity of a node in the Hanzo decentralized AI
cloud: how the identity is created, how it is held, what it signs, what a
verifier may conclude from what it signs, how it stakes, how it is paid, and how
it ends.

It does not specify consensus, the scheduler, the workload API, or the
verification of AI work. Verification of advertised capability and delivered
service is HIP-0095 (QoS challenges) and HIP-0096 (compute rewards); this
document states only where the identity boundary meets those and what identity
does not supply.

Throughout, \emph{observed} marks a statement about code or hardware that was
read or measured, with the location given. \emph{Proposed} marks a change this
document asks for. Statements with neither marker are definitions or
consequences of the two.

The key words MUST, MUST NOT, SHOULD, SHOULD NOT and MAY are to be interpreted as
in RFC 2119.

% =====================================================================
\section{Terms}

\begin{description}
\item[Node] A host that runs \code{hanzod} and offers compute to the network.
  Physical, and heterogeneous: Appendix~\ref{app:fleet} records three, on three
  architectures, with three different attestation capabilities.

\item[Machine] A microVM the node starts. A node serves many. The vocabulary
  \emph{machine} $>$ \emph{cluster} $>$ \emph{vm} $>$ \emph{sandbox} is fixed
  elsewhere in the platform; in this document \emph{machine} means a microVM the
  node runs, and \emph{host} means the physical box.

\item[Monitor] The \code{hanzo-vm} binary that builds and runs a machine. VZ on
  macOS, cloud-hypervisor or KVM on Linux.

\item[Node microVM] The one machine on each host that joins the k3s cluster and
  runs the Hanzo cloud. It is the uniformity layer: it is why an
  Apple-Silicon Mac and an x86 workstation present the same surface to the
  cluster. Exactly one per node.

\item[Boot] The set of inputs that determine a machine: monitor binary, OS-asset
  version, kernel, root filesystem, shape, and argv. Observed as \code{Boot} in
  \code{hanzo-attest/src/boot.rs}.

\item[Measurement] $\mu$, the 32-byte BLAKE3 digest over the canonical encoding
  of a \code{Boot}. Observed; the encoding is reproduced in
  Appendix~\ref{app:encoding}.

\item[Node key] The secp256k1 key whose address is the node's identity. Held by
  \code{hanzod}. Signs the node's statements about itself.

\item[Node address] $A_n = \addr(pk_n)$, the last 20 bytes of
  $\mathrm{Keccak256}$ of the uncompressed public key. This is the identity.

\item[Owner key] The key that registers the node, posts the bond, ends the
  validation, and withdraws. Address $A_o$. Held by the operator, off the node.

\item[Evidence] A signed statement about a measurement. Observed as
  \code{Evidence} in \code{hanzo-attest/src/evidence.rs}.

\item[Assurance class] What a node can \emph{prove} about the platform holding
  its key. Derived from evidence, never declared. \S\ref{sec:class}.
\end{description}

% =====================================================================
\section{What signatures, measurements and attestations establish}
\label{sec:proves}

This section is the foundation. Everything later --- staking, slashing, reward
attribution, fleet enumeration --- inherits exactly the assurance established
here and no more.

Let $\Sigma = (\mathsf{Gen}, \Sign, \Vrfy)$ be a signature scheme existentially
unforgeable under chosen-message attack. Fix a keypair $(sk, pk)$ and write
$A = \addr(pk)$. Over a run of the system define three predicates:

\begin{description}
\item[$\Held(P, sk, t)$] party $P$ can evaluate $\Sign_{sk}(\cdot)$ at time $t$.
\item[$\Ran(h, I, t)$] host $h$ was executing image $I$ at time $t$.
\item[$\Conf(sk, h)$] $sk$ is confined to $h$: no party other than $h$ can
  evaluate $\Sign_{sk}(\cdot)$.
\end{description}

\subsection{A signature}

\begin{claim}[What verification gives]\label{cl:sig}
If $\Vrfy_{pk}(m, \sigma) = 1$ then, under EUF-CMA, there exist a party $P$ and a
time $t$ with $\Held(P, sk, t)$, and $P$ chose the message $m$.
\end{claim}

\begin{claim}[What it does not give]\label{cl:indep}
The conclusion of Claim~\ref{cl:sig} is independent of $\Ran$ and $\Conf$. For
any host $h$ and image $I$ a verifier believes in, there is a run in which the
conclusion holds and $\Ran(h, I, t)$ is false.
\end{claim}

\noindent The proof is the attack: copy $sk$ to any machine and sign $m$ there.
Nothing in $\sigma$ distinguishes the two runs, because nothing in $\sigma$
depends on where the signing happened.

Claim~\ref{cl:indep} holds for every message, including a message that is
\emph{about} a machine. Signing a boot measurement does not make the signature
evidence about a boot. It makes it evidence that the key-holder asserted a
number.

\subsection{A measurement}

$\mu$ is a function of files on a host, computed by the host, before the machine
starts. Two parties holding the same inputs compute the same $\mu$; a third party
holding only $\mu$ can check a claim about those inputs. That is its entire
content.

\begin{claim}\label{cl:meas}
$\mu$ binds inputs. It does not witness execution. A host that computes $\mu$
over one set of files and hands the monitor a different set produces a
measurement that verifies against the first set and describes nothing that ran.
\end{claim}

This is stated in the implementation itself
(\code{hanzo-attest/src/boot.rs}, module documentation: ``The measurement is
taken by the host, over files on the host, before the machine starts. It binds
the inputs; it does not witness the execution.''). We restate it because it is
the property most often lost in transit: a measurement in a registry looks like
a fact about a running system and is not one.

\subsection{An attestation}
\label{sec:att}

\begin{claim}[What a hardware quote gives]\label{cl:att}
Let $Q$ be a quote over $(\mu, n, d)$ verifying under vendor root $R$, where $n$
is a nonce the verifier chose. Then there exists a platform $h$ whose attestation
key $R$ certifies, such that $h$'s measured state was $\mu$ at a time after $n$
was issued, and $h$ reported the value $d$.
\end{claim}

Three things do not follow:

\begin{enumerate}[label=(\alph*)]
\item \textbf{$\mu$ is not behavior.} A measurement covers what was loaded. Code
  loaded later, a kernel module, a \code{kexec}, a writable root filesystem, or
  an interpreter reading a script from the network all change behavior without
  changing $\mu$. A measured platform running a program that fetches and executes
  arbitrary bytes measures as honestly as one that does not.

\item \textbf{$d$ is not $\Conf$.} Reporting $d = H(pk_n)$ inside a quote binds
  the key to the platform's \emph{statement}, not to the platform. $\Conf(sk_n,
  h)$ follows only if $sk_n$ was generated inside $h$ and is non-extractable, and
  only if the platform reports that fact in a form the verifier checks. A key
  generated elsewhere and copied in attests exactly as well.

\item \textbf{The host is in the TCB unless the VM is confidential.} Guest-side
  measurement of a machine whose memory the hypervisor can read and write
  establishes $\Ran$ only relative to an honest host. On AMD SEV-SNP, Intel TDX,
  or Arm CCA the host is removed from that base. \emph{Observed:} none of the
  three hosts in Appendix~\ref{app:fleet} runs a confidential VM today ---
  \code{kvm\_amd.sev\_snp = N} on the x86 host, no Realm support on the aarch64
  host, and Virtualization.framework on the Mac. Guest measurement on this fleet
  therefore defends against a remote attacker and not against the node operator.
\end{enumerate}

Point (c) matters most for third-party nodes, where the node operator is the
adversary of interest. On the current fleet the only hardware root that speaks
about the \emph{host} is the TPM (\S\ref{sec:class}).

\subsection{The possession--platform gap}
\label{sec:gap}

\begin{definition}[Possession--platform gap]
The gap between $\Held$ and $\Conf$: a signature establishes that \emph{someone}
could sign, and never that \emph{this platform} is the only one that could.
\end{definition}

Every downstream property inherits it. If the node key can be copied, then
``which node produced this statement'' has no answer, and with it neither does
``which node earned this reward'', ``which node's stake answers for this fault'',
nor ``how many nodes are in the fleet''. Fleet enumeration by wallet is only as
meaningful as the injectivity of the map from wallets to machines, and a copied
key breaks injectivity silently: two machines answer as one node, and the console
shows one row.

Closing the gap requires two properties, both of which must be reported and
checked, not assumed:

\begin{enumerate}[label=(\roman*)]
\item the key is generated inside a platform and is non-extractable from it, and
\item the platform can produce, under a root the verifier trusts, a statement
  binding that key to a measured state.
\end{enumerate}

No software arrangement supplies (i). Custody in a key-management service does
not supply it either --- it \emph{moves} the gap, which is a real improvement and
a different property; see \S\ref{sec:ladder}.

\subsection{What an unattested but signed join does not prove}
\label{sec:notprove}

A node that signs its boot measurement with its node key and presents that to the
network --- \code{Kind::Host} in the implemented code --- proves that the holder
of $sk_n$ asserted $\mu$ at a stated time and that the assertion has not been
altered since. It does not prove any of the following. Each is a distinct failure,
and none is covered by adding more signatures of the same kind:

\begin{enumerate}[label=\textbf{N\arabic*.},leftmargin=2.6em]
\item \textbf{Not that the machine ran the measured image.} The measurement is
  computed over files chosen by the same party that signs. A node that measures
  image $I$ and boots image $I'$ produces evidence that verifies
  (Claim~\ref{cl:meas}).
\item \textbf{Not that the statement came from the host it describes.} The
  signing key is portable; the statement can be produced anywhere, including on
  a laptop, and relayed.
\item \textbf{Not that the key is on any particular host at all.} $\Conf$ is not
  established (Claim~\ref{cl:indep}).
\item \textbf{Not that the identity is one machine.} Without an epoch-unique
  statement (\S\ref{sec:equiv}), $N$ machines holding a copy of $sk_n$ join as
  one node and are indistinguishable from one machine.
\item \textbf{Not that the advertised inventory is real.} CPU, memory, and GPU
  counts are numbers the node reports. The \code{Shape} inside the measurement is
  what \code{hanzod} \emph{asked the monitor for}, which is what it can honestly
  sell, and is not a statement about the physical host.
\item \textbf{Not that the machine still runs what it reported.} The measurement
  is taken once, before start. Nothing in the statement constrains the machine
  afterwards.
\item \textbf{Not that the node will do the work correctly.} Identity is
  orthogonal to service. HIP-0095 and HIP-0096 own that.
\item \textbf{Not that the signer is authorized.} Verification of a signature is
  not authorization. Whether the network accepts $A_n$ is a registration question
  (\S\ref{sec:stake}), and the implementation is right to separate them
  (\code{hanzo-attest/src/evidence.rs}: ``Deciding that a signer is acceptable is
  a policy question \dots\ not here'').
\end{enumerate}

\subsection{Five controls, none substitutable}

\begin{center}
\small
\begin{tabular}{@{}>{\raggedright\arraybackslash}p{0.15\textwidth}>{\raggedright\arraybackslash}p{0.35\textwidth}>{\raggedright\arraybackslash}p{0.40\textwidth}@{}}
\toprule
\textbf{Control} & \textbf{Establishes} & \textbf{Does not establish} \\
\midrule
Signature &
Authorship of a statement under a key &
Which machine; which image; uniqueness \\[3pt]
Custody &
Who can cause a signature to exist &
That the statement is true \\[3pt]
Attestation &
That a vendor-rooted platform reported a measured state, bound to a fresh nonce &
That the state is the behavior; that the host is outside the TCB, unless the VM is confidential \\[3pt]
Stake &
A price for a fault, once the fault is established &
Detection \\[3pt]
Challenge (HIP-0095) &
Detection of false capability and service claims &
Prevention; identity \\
\bottomrule
\end{tabular}
\end{center}

Layering is not redundancy here: each control's ``does not establish'' column is
covered by a different row, and removing any row leaves a hole no other row
fills. A system with attestation and no stake detects faults it cannot price. A
system with stake and no attestation prices faults it cannot detect.

% =====================================================================
\section{The identity}
\label{sec:identity}

\subsection{Two roles, one identity}
\label{sec:roles}

A single key cannot do both of the jobs a node needs done. The node key signs
constantly, from a box that runs other people's workloads and is, for
third-party nodes, outside our physical control. The bond must be movable only
by someone we still trust after that box is compromised. If they are the same
key, compromise of the node key also lets the attacker withdraw the bond, so the
economic control intended to price the compromise is unavailable in the case it
was posted against.

Therefore:

\begin{description}
\item[Node key ($A_n$)] Lives with \code{hanzod}. Signs epoch statements and
  machine records. \emph{Is the node identity.} Has no authority to move value.
\item[Owner key ($A_o$)] Lives with the operator, off the node. Registers the
  node, posts and withdraws the bond, ends the validation, changes the reward
  recipient. Never on the box.
\end{description}

This is one identity with an explicit controller, not two identities. $A_n$ is
what the fleet is enumerated by, what evidence is signed with, what work is
attributed to, and what rewards are paid to. $A_o$ appears once per lifecycle
event.

\emph{Observed:} the vendored validator manager sets
\code{\_posValidatorInfo[validationID].owner = \_msgSender()} at registration
(\code{luxfi/standard}, \code{contracts/validators/PoSValidatorManager.sol}:440--442)
and authorizes end-of-validation, reward withdrawal, and recipient changes against
that one address. There is no second address field. The registration therefore
knows $A_o$ and does not know $A_n$; supplying the binding is
\hyperref[sec:req]{IR-3}.

\begin{remark}
For a Hanzo-operated node, $A_o$ SHOULD be a \code{KindTreasury} or
\code{KindMPC} wallet --- threshold custody, no single holder --- and MUST NOT be
a key present on any node. \emph{Observed:} both custody kinds exist in
\code{hanzoai/cloud}, \code{apps/wallet/custody.go}, and a
\code{TierValidator} wallet tier already exists in the same file.
\end{remark}

\subsection{Derivation}

$A_n$ is an EVM-runtime account address: the low 20 bytes of
$\mathrm{Keccak256}$ over the uncompressed secp256k1 public key. \emph{Observed}
in \code{hanzoai/cli}, \code{src/commands/wallet.rs} (\code{address\_from\_vk}),
with the canonical path
\code{m/44\textquotesingle/60\textquotesingle/0\textquotesingle/0/0} for locally
derived keys.

The node key MUST be generated for the node and MUST NOT be derived from an
operator seed that also derives other nodes' keys. Derivation from a shared
mnemonic makes every node's key recoverable from one secret, which converts a
single custody failure into a fleet-wide one. Where a fleet seed is used for
operational reasons, each node's key MUST come from a distinct hardened account
index and the seed MUST NOT be present on any node.

\subsection{The one object}

A node record is:
\[
  \bigl(A_n,\; A_o,\; \text{chain},\; \mu_{\text{node}},\; \text{class},\;
  \text{validationID},\; \text{status}\bigr)
\]
and nothing else identifies a node anywhere in the system. \code{admin.hanzo.ai}
enumerates by $A_n$. The scheduler selects by $A_n$ and \code{class}. Settlement
pays $A_n$. The registry keys on $A_n$.

\emph{Observed gap:} \code{hanzo-attest}'s \code{Record} carries
\code{node: String} --- a free-form name --- and its \code{Signer} is an ed25519
public key (\code{evidence.rs}). The node's identity secret today is an ed25519
key held as a hex string in a plaintext \code{KEY=value} file. In
\code{hanzoai/node}, \code{hanzo-bin/hanzo-node/src/utils/keys.rs},
\code{generate\_or\_load\_keys} reads \code{IDENTITY\_SECRET\_KEY} from that file,
falls back to the environment variable of the same name, and otherwise generates
a fresh key; \code{src/runner.rs}:201--212 writes the result back with
\code{std::fs::write}, which creates at the process umask rather than
owner-only. (A comment at \code{keys.rs}:72 refers to Stronghold; no Stronghold
path exists in that file.) That is L0 custody by the ladder of
\S\ref{sec:ladder}: a copyable key, and therefore a forgeable identity.
Replacing it is \hyperref[sec:req]{IR-1}.

\emph{Relationship to HIP-0097.} HIP-0097 specifies \code{did:hanzo:<name>} as
the DID method for every Hanzo subject. $A_n$ is the node's key material and its
settlement address; a \code{did:hanzo:} name MAY resolve to a DID Document whose
verification method is $A_n$. Two names for one subject is acceptable only in
that direction --- the DID resolves to the address; the address is not derived
from the DID --- because settlement and staking cannot depend on a resolver being
reachable.

% =====================================================================
\section{Custody}
\label{sec:custody}

\subsection{The ladder}
\label{sec:ladder}

Custody options differ in one measurable respect: what an attacker must obtain in
order to sign as $A_n$.

\begin{center}
\small
\begin{tabular}{@{}l>{\raggedright\arraybackslash}p{0.30\textwidth}>{\raggedright\arraybackslash}p{0.44\textwidth}@{}}
\toprule
& \textbf{Custody} & \textbf{To forge the identity, obtain} \\
\midrule
L0 & Plaintext file or environment variable & File read as any user that can read it. \emph{This is the current state.} \\
L1 & OS keychain, owner-only file & Code execution as the node's user. \\
L2 & Held by KMS/MPC; node never sees key & A credential that authorizes KMS to sign --- which is a bearer token on the node. \\
L3 & Sealed to a platform: generated inside, non-extractable, released only under a policy over measured state & Physical compromise of the platform, or a break in its root of trust. \\
\bottomrule
\end{tabular}
\end{center}

L2 is where the interesting mistake lives. Moving the key into KMS is a genuine
improvement --- the key is not on the box, so a stolen disk is not a stolen
identity, and every signature is centrally logged and revocable. But it does not
close the possession--platform gap; it \emph{relocates} it. The question ``which
machine holds the key'' becomes ``which principal can make KMS sign'', and on a
node that principal is authenticated by a credential that sits on the node. Under
L2 the identity is exactly as strong as the custody of that credential --- which,
absent hardware, is L0 or L1 again.

\begin{claim}[The recursion]\label{cl:recursion}
Software custody delegation does not increase assurance about \emph{location}. It
replaces custody of a signing key with custody of an authorization credential.
The chain terminates only at a key that some hardware refuses to export.
\end{claim}

This is not an argument against KMS. It is an argument for what KMS custody must
be combined with, and it produces the requirement below.

\subsection{Requirement}

\begin{enumerate}[label=\textbf{C\arabic*.},leftmargin=2.6em]
\item The node key MUST NOT exist in plaintext at rest anywhere: not in an
  environment variable, not in a file, not in a container image, not in a
  Kubernetes Secret that is not itself KMS-sealed.
\item The node key SHOULD be held at L2 --- generated inside KMS/MPC, never
  returned --- with \code{hanzod} holding only $A_n$ and a signing capability.
  \emph{Observed:} this is exactly the cloud-custody model already implemented in
  \code{hanzo wallet} (``keys are derived and held server-side \dots\ and NEVER
  leave it --- the CLI only ever sees the address'') and in
  \code{cloud/apps/wallet/custody.go} (\code{KindKMS}, \code{KindMPC}).
\item Where the node key is at L2, the credential that authorizes signing MUST be
  scoped to (a) this node's wallet only, (b) the epoch-statement and
  machine-record message shapes only, and (c) a bounded rate. A node's signing
  credential MUST NOT be able to sign an arbitrary transaction, because a
  credential that can sign anything is the owner key by another name.
\item On a platform with a TPM 2.0, the credential in C3 MUST be sealed to a PCR
  policy, so that a credential lifted from the disk is unusable on a different
  platform or after a boot-chain change.
\item On a platform with a hardware enclave (SEV-SNP, TDX, CCA), the key SHOULD
  be generated inside the enclave and its release policy bound to the enclave
  measurement, reaching L3.
\item Secrets MUST be stored only in Hanzo KMS. Where a digest suffices, the
  digest is stored and the value is not.
\end{enumerate}

\emph{Observed} for the KMS envelope this relies on: secrets are sealed as
AES-256-GCM over a per-secret DEK with the DEK wrapped by ML-KEM
(\code{luxfi/kms}, \code{pkg/store/secrets.go}, \code{ModeStandard = "aead+mlkem"}),
the AAD binds path, name, environment and version, and
\code{pkg/store/crypto\_test.go:TestAADBindingPreventsSwap} demonstrates that
swapping one secret's wrapped DEK into another fails. Transport is a ZAP
handshake combining X25519 and ML-KEM-768 under HKDF-SHA256
(\code{pkg/zap/handshake.go}).

\subsection{Measurement-gated release}

At L3 the release policy is the point of the exercise, so it is specified rather
than left to configuration:

\begin{enumerate}[label=\textbf{G\arabic*.},leftmargin=2.6em]
\item The key material is generated inside the platform and is not exportable.
\item A quote over $(\mu_{\text{node}}, n, H(pk_n))$, with $n$ chosen by the
  verifier, is the only evidence that binds the key to the platform. It MUST be
  requested with a fresh $n$ per verification; a stored quote is a statement
  about the past.
\item Release --- or use --- MUST be denied when the measured state is not on the
  allowed list for the node's declared image. Denial MUST be the default on any
  error, including an unreachable verifier.
\item The allowed list is a set of $\mu$ values, published, and changed by a
  release of the image --- not by an operator edit on the node.
\end{enumerate}

\subsection{Assurance class}
\label{sec:class}

A node's assurance class is derived from evidence the network verified. It is a
field on the node record, it is visible in the console and to the scheduler, and
a node MUST NOT be able to set it.

\begin{center}
\small
\begin{tabular}{@{}l>{\raggedright\arraybackslash}p{0.22\textwidth}>{\raggedright\arraybackslash}p{0.50\textwidth}@{}}
\toprule
\textbf{Class} & \textbf{Root} & \textbf{Requires} \\
\midrule
\code{self} & the node's own key &
An epoch statement verifying under $A_n$. Trust in it is trust in the node
operator. \emph{Implemented} (\code{Kind::Host}). \\[3pt]
\code{platform} & TPM 2.0 &
A quote over PCRs binding $H(pk_n)$ to a fresh nonce, \emph{and} an enforced
verified boot chain. Without enforcement the PCR values are attacker-chosen and
the class is \code{self}. \emph{Proposed.} \\[3pt]
\code{confidential} & SEV-SNP, TDX, CCA &
A vendor-rooted report; removes the host from the TCB. \emph{Proposed; no
hardware on the current fleet.} \\
\bottomrule
\end{tabular}
\end{center}

\emph{Observed} on the current fleet (Appendix~\ref{app:fleet}): the aarch64 host
has TPM 2.0 with a SHA-256 PCR bank, an event log at
\code{/sys/kernel/security/tpm0/binary\_bios\_measurements}, and SecureBoot
enabled --- it can reach \code{platform}. The x86 host has TPM 2.0 with SecureBoot
\emph{disabled} and \code{sev\_snp = N} --- it is \code{self} until SecureBoot is
enabled, and cannot reach \code{confidential} on the running configuration. The
Apple-Silicon host exposes no TPM and no confidential-VM facility to a Linux
guest --- it is \code{self}, and this is a property of the platform, not a
deficiency to be worked around.

\begin{remark}[Why the Mac is in the fleet anyway]
The microVM is the uniformity layer: it is what lets a Mac present the same k3s
node as a Linux box. Uniformity of \emph{interface} is not uniformity of
\emph{assurance}, and the design does not pretend otherwise. The Mac is a first-class
node at class \code{self}, and the scheduler is what decides which workloads a
\code{self} node may receive. One identity primitive, one enumerable field for
the difference.
\end{remark}

The existing implementation names \code{Kind::Snp} and \code{Kind::Tdx} and
implements neither. Given the hardware actually present, \emph{a TPM 2.0 attester
and verifier is worth more than either}, because two of three hosts have the
hardware today and zero have SNP or TDX. That is \hyperref[sec:req]{IR-4}.

% =====================================================================
\section{Joining}
\label{sec:join}

\subsection{The epoch statement}

Let $T$ be the epoch length and $e = \lfloor t/T \rfloor$. A node publishes one
\emph{epoch statement} per epoch:
\[
S = \bigl(\text{tag},\; \text{chainID},\; A_n,\; e,\; \mu_{\text{node}},\;
\text{class evidence}\bigr),\quad \sigma = \Sign_{sk_n}(S).
\]

Four fields carry weight and none is decorative:

\begin{description}
\item[chainID] Prevents a statement made for testnet from being replayed on
  mainnet. \emph{Observed:} the sovereign L1 uses $\text{networkID} =
  \text{chainID}$, mainnet 36963, testnet 36962, devnet 36964
  (\code{hanzoai/cli}, \code{src/commands/network.rs}).
\item[$A_n$] Binds the statement to the identity it is filed under. Without it, a
  statement can be presented under a different registration.
\item[$e$] Makes the statement unique per epoch, which is what makes equivocation
  provable (\S\ref{sec:equiv}) and what supplies liveness without a second
  heartbeat protocol.
\item[$\mu_{\text{node}}$] The measurement of the node microVM --- the one that
  joins the cluster --- not of a tenant machine.
\end{description}

The epoch statement is distinct from the per-machine \code{Record} already
implemented. A node serves many machines and files many records; it has exactly
one node microVM and files exactly one epoch statement per epoch. Conflating them
would make honest multi-tenancy look like equivocation.

\emph{Observed gap:} the implemented preimage is
$(\text{tag}, \text{kind}, \mu, \text{timestamp})$
(\code{hanzo-attest/src/evidence.rs}, \code{preimage}). It carries no chain, no
address, and no epoch. It is therefore presentable under any registration on any
chain, and equivocation over it is not defined. Adding the three fields is
\hyperref[sec:req]{IR-2}, and it is the single highest-value change in this
document.

\subsection{What the network verifies}

On receiving a join or an epoch statement, the verifier MUST, in order, and
refusing at the first failure:

\begin{enumerate}[label=\textbf{V\arabic*.},leftmargin=2.6em]
\item \textbf{Shape.} $S$ parses; \code{chainID} is this chain; $e$ is the
  current epoch or the immediately previous one.
\item \textbf{Signature.} $\Vrfy_{A_n}(S, \sigma) = 1$.
\item \textbf{Coherence.} $\mu_{\text{node}}$ equals the measurement recomputed
  from the accompanying \code{Boot}. \emph{Observed:} \code{Record::coherent}
  already performs exactly this check.
\item \textbf{Registration.} $A_n$ is bound to an active \code{validationID}
  whose bond meets the minimum, and the binding was written by $A_o$.
\item \textbf{Uniqueness.} No other statement exists for $(A_n, e)$ with a
  different $\mu_{\text{node}}$. A second, differing statement is not an error
  to discard --- it is evidence to keep (\S\ref{sec:equiv}).
\item \textbf{Image.} $\mu_{\text{node}}$ is on the allowed list for the release
  the node claims. A node running an unpublished image is admitted at class
  \code{self} with no image assertion, or refused, according to network policy;
  it MUST NOT be admitted as though it were running a published image.
\item \textbf{Class.} Any hardware evidence verifies under the vendor root, is
  bound to a nonce this verifier issued within the freshness window, and reports
  $H(pk_n)$. The class recorded is the highest class \emph{proved}, defaulting to
  \code{self}.
\end{enumerate}

Every failure is a refusal. A verifier that cannot reach the vendor root records
class \code{self}; it MUST NOT record a higher class on an unverified claim, and
it MUST NOT refuse the join for that reason alone --- the two are different
decisions and conflating them turns an availability failure into a fleet outage.

\subsection{What the network has after V1--V7}

Precisely: that the holder of $sk_n$, bound to an active bonded registration,
asserted image $\mu_{\text{node}}$ for epoch $e$ on this chain, and that at class
\code{platform} or above a vendor-rooted platform corroborated the state and the
key binding. Nothing in \S\ref{sec:notprove} that is not covered by class
\code{platform} or above becomes true because V1--V7 passed.

% =====================================================================
\section{Stake and forfeiture}
\label{sec:stake}

\subsection{What exists}

\emph{Observed}, in \code{luxfi/standard},
\code{contracts/validators/} --- the vendored ACP-77 validator manager, renamed
to Network/L1 nomenclature with the ABI preserved:

\begin{itemize}
\item \code{initializeValidatorRegistration} is \code{payable}. The native value
  sent stays in the manager and is the bond; \code{\_lock} records the amount and
  \code{\_unlock} returns it with \code{payable(to).sendValue(value)}
  (\code{NativeTokenStakingManager.sol}:68--94).
\item The registration is keyed by \code{validationID} (\code{bytes32}) and
  records \code{owner = \_msgSender()}, a delegation fee, a minimum stake
  duration, and an uptime counter (\code{PoSValidatorManager.sol}:440--445).
\item \code{submitUptimeProof} and the end-of-validation path condition rewards
  on uptime; below the threshold, rewards are forfeit and the bond is still
  returned.
\item The registration input carries \code{nodeID} (\code{bytes}) and
  \code{blsPublicKey}. The contract checks the BLS key's \emph{length} (48) and
  that \code{nodeID} is non-empty and not already registered
  (\code{ValidatorManager.sol}:249--281). It does not verify possession of
  either key.
\item \textbf{There is no slashing.} A case-insensitive search for
  \code{slash} across \code{contracts/validators/} returns nothing.
\end{itemize}

So the economic instrument that exists today is uptime-conditioned reward
forfeiture plus bond return. Any statement that a bad attestation ``slashes the
stake'' is, as of this writing, a proposal. This document treats it as one.

\subsection{Registration binds the identity}

\emph{Proposed.} A separate, minimal registry contract --- not a fork of the
vendored manager, whose ABI is preserved deliberately for in-place upgradability:

\begin{quote}\ttfamily
mapping(bytes32 validationID => address node) nodes;\\
mapping(address node => bytes32 validationID) registrations;\\
\\
function bind(bytes32 validationID, address node) external;\\
event Bound(bytes32 indexed validationID, address indexed node);
\end{quote}

\code{bind} is callable only by the \code{owner} of \code{validationID} as the
manager reports it, refuses a node address already bound to a live registration,
and refuses to rebind a live registration to a different node address. One
mapping pair, one function, one event. Everything else --- weight, churn, uptime,
delegation --- stays where it already works.

This is the minimum that makes $A_n$ mean something on-chain, and it is what lets
the scheduler, the settlement path, and the console agree on one identifier
without any of them trusting the others' view.

\subsection{Forfeiture for self-proving faults}
\label{sec:forfeit}

\emph{Proposed.} On-chain forfeiture applies to faults that are
\emph{self-proving}: a verifier can establish them from the submitted evidence
alone, in constant time, without a vote, an oracle, or a trusted reporter. This
restriction is the whole design, and it is what keeps the mechanism from becoming
a governance surface.

\begin{quote}\ttfamily
function forfeit(bytes32 validationID, bytes calldata proof) external;\\
event Forfeited(bytes32 indexed validationID, uint256 amount, bytes32 fault);
\end{quote}

The contract verifies \code{proof} itself and burns or redistributes the bond
according to the network's reward policy. Anyone may submit; a submitter reward
paid from the forfeited amount is what makes detection someone's job.

Faults that are \emph{not} self-proving --- poor service, wrong answers,
unavailability, a false inventory claim --- are handled by the mechanisms that
already exist for them: uptime-conditioned reward forfeiture in the manager, QoS
challenges (HIP-0095), and delisting by the scheduler. They MUST NOT be routed
through \code{forfeit}, because each would require a trusted reporter, and a
trusted reporter that can burn a stranger's bond is a larger risk than the
behavior it polices.

\subsection{Equivocation}
\label{sec:equiv}

\begin{definition}[Equivocation]
Two epoch statements $S_1, S_2$ with the same $(\text{chainID}, A_n, e)$,
$\mu_1 \neq \mu_2$, both verifying under $A_n$.
\end{definition}

\begin{claim}
Equivocation is publicly verifiable in two signature verifications and one
inequality, and an honest node running one node microVM cannot produce it.
\end{claim}

Equivocation is what a copied key produces on-chain. If $sk_n$ exists on two
machines and both join, the two report different boots in the same epoch and the
bond is forfeit. This is what makes the custody requirement enforceable rather
than advisory: an operator who copies the node key across machines does not face
a probability of detection, because detection requires no monitoring --- the two
statements are the proof, and either party to the collision can submit them.

The rule also gives the honest operator a defined failure mode. A node restarted with a
new image mid-epoch produces a second statement with a different $\mu$ and would
equivocate against itself. The rule is therefore that a node MUST NOT publish a
second statement in an epoch it has already published in; a restart waits for the
next epoch, or publishes an explicit \emph{supersede} that names the prior
statement's digest and is signed by the same key --- which is a chain, not a
fork, and is not equivocation. $T$ MUST be short enough that waiting is cheap
(minutes, not hours) and long enough that ordinary restarts do not dominate the
statement volume.

% =====================================================================
\section{Reward}

Rewards pay $A_n$ --- the identity address --- from two independent sources:

\begin{description}
\item[Validation rewards] From the validator manager, conditioned on uptime.
  \emph{Observed:} the manager keeps a \code{\_rewardRecipients} mapping that
  \emph{defaults to} \code{owner} (\code{PoSValidatorManager.sol}:330) and can be
  pointed elsewhere by \code{changeValidatorRewardRecipient}
  (\code{PoSValidatorManager.sol}:254--266).
\item[Work rewards] From compute settlement (HIP-0096), for AI work the node
  actually served.
\end{description}

Two consequences follow from the observed divergence and both must be stated
rather than assumed away:

\begin{enumerate}[label=\textbf{R\arabic*.},leftmargin=2.6em]
\item The reward recipient MUST default to the identity address. Divergence is
  permitted --- an operator consolidating payouts is a legitimate need --- but it
  is an explicit on-chain act, visible in the fleet view next to the node it
  applies to.
\item Accounting, attribution, and enumeration key on $A_n$, \textbf{never} on
  the recipient. A recipient is where money goes; it is not who did the work, it
  is not unique to a node, and treating it as an identifier silently merges
  distinct nodes into one row.
\end{enumerate}

Paying the address the node signs with means ``which node earned this'' and
``which node said this'' resolve through the same lookup. That is the practical
content of using one primitive for attestation, staking, payment and enumeration
instead of four identifiers that have to be reconciled. R2 is what keeps it true
once a recipient diverges: divergence is permitted, and it must not propagate
into the identifier.

% =====================================================================
\section{Lifecycle}
\label{sec:lifecycle}

\subsection{Rotation is succession}

An EVM address is a hash of the public key. There is no operation that changes
the key and keeps the address. A design that wanted stable identifiers across key
changes would need a level of indirection, and indirection here would be a
mistake: after a compromise, carrying the identifier forward carries the
compromised node's standing forward with it.

So rotation is \textbf{succession}: a new key, a new address, an explicit link.

\begin{enumerate}[label=\textbf{S\arabic*.},leftmargin=2.6em]
\item The operator generates $sk_n'$ at the target assurance class and obtains
  $A_n'$.
\item $A_o$ calls \code{bind} to move the registration's node address from $A_n$
  to $A_n'$, emitting \code{Bound}. The bond does not move; the registration does
  not end; no unbonding period is incurred.
\item The outgoing key SHOULD sign a statement naming $A_n'$, which lets an
  observer follow the chain without consulting the registry. It is corroboration,
  not authority: the authority for succession is $A_o$, precisely because the
  case that matters most is the one where $sk_n$ is in someone else's hands.
\item Assurance class does not transfer. The successor is measured on what it can
  prove.
\item Accrued reputation and service history transfer only when the succession is
  \emph{planned} --- a hardware refresh or a key-hygiene rotation. After a
  suspected compromise they MUST NOT transfer.
\end{enumerate}

Because $A_o$ is the authority, succession works when the node key is
compromised, which is the case it exists for. This is the concrete reason for the
two-key split in \S\ref{sec:roles}: with one key there is no step S2 an honest
operator can take that the attacker cannot.

\subsection{Revocation}

Revocation is immediate and off-chain-effective:

\begin{itemize}
\item $A_o$, or the network on a proven fault, marks $A_n$ revoked.
\item The scheduler stops placing work on the node at once. Running work is
  drained under the workload's own policy.
\item Statements signed by $A_n$ after the revocation timestamp are refused.
\item The bond does \textbf{not} unlock at revocation. It unlocks after the
  unbonding period, so that a node cannot escape a forfeiture that is still being
  submitted. Making revocation fast and unbonding slow is the whole point of
  separating them.
\end{itemize}

\subsection{Decommission}

The orderly path, for a node leaving without fault: drain, stop accepting work,
publish a final epoch statement marked terminal, and let $A_o$ call the manager's
existing end-of-validation path. Uptime is settled, rewards pay out, the bond
returns to $A_o$ after the unbonding period, and the record is retired.

$A_n$ MUST NOT be reused. A retired address has an unknowable custody history:
whoever held the key still holds it, and nothing distinguishes a returning node
from an old copy.

\subsection{Compromise}

There is no procedure that un-leaks a key. What the procedure can do is stop the
attacker from profiting and stop the network from believing the wrong things.
Assume $sk_n$ is disclosed at unknown time $t_c$:

\begin{enumerate}[label=\textbf{K\arabic*.},leftmargin=2.6em]
\item \textbf{Revoke immediately} (previous section). Scheduling stops; the bond
  stays locked.
\item \textbf{Do not use the node key.} It is the attacker's key too. Every
  recovery action goes through $A_o$.
\item \textbf{Treat statements as unattributable from $t_c$.} Not false ---
  unattributable. Work attributed to $A_n$ after $t_c$ cannot be assigned to the
  operator or the attacker on the evidence available.
\item \textbf{Expect equivocation, and let it settle.} If the attacker runs the
  key, the two statements collide and the bond is forfeit. That is the mechanism
  working: the bond is what the operator posted against exactly this.
\item \textbf{Succeed to $A_n'$} at a class at least as high as the one that
  failed. If the failure was custody, the successor MUST be at a strictly higher
  class, or the same failure is available again.
\item \textbf{Rotate the credential, not only the key.} At L2 the compromise may
  be of the signing \emph{authorization}, not the key --- see
  Claim~\ref{cl:recursion}. Revoke the node's KMS credential, and check the KMS
  audit log for signatures in $[t_c, \text{now}]$ that the node cannot account
  for.
\end{enumerate}

\begin{remark}
If $A_o$ is compromised, this is not a node incident. The bond is movable and the
registration is controllable by the attacker. Recovery is at the operator level:
threshold custody for $A_o$ is what makes the event survivable, which is why
\S\ref{sec:roles} requires it rather than suggesting it.
\end{remark}

% =====================================================================
\section{Hanzo-operated and third-party nodes}
\label{sec:same}

They are the same object. A Hanzo node is a node whose owner key Hanzo holds.
There is no Hanzo-only field, no privileged registration path, no separate
enumeration, and no code path that branches on who operates the box. The console
lists them in one table; the scheduler selects from one pool; settlement pays
them through one mechanism.

Three differences remain. Each is a value on the record, and none is a difference
in the identity:

\begin{enumerate}[label=\textbf{D\arabic*.},leftmargin=2.6em]
\item \textbf{Who can authorize signing.} For a Hanzo node, $A_o$ and the node's
  KMS credential live in Hanzo's KMS under Hanzo's IAM org; for a third party,
  in theirs. This changes who can cause a signature, not what the signature
  means. Both are subject to \S\ref{sec:custody} identically.
\item \textbf{Physical control, and therefore who the adversary is.} For a
  third-party node, the operator is a possible adversary, and
  \S\ref{sec:att}, point~(c), applies at full strength: without a confidential VM,
  guest measurement does not defend against the operator. For a Hanzo node the
  host is under our control, which is a real assurance difference and is
  \emph{not} expressible as a cryptographic proof to a third party. It is
  recorded as the operator org on the node record and priced by policy, not
  asserted as a security property.
\item \textbf{Whose money is at risk.} Hanzo's own nodes are bonded by Hanzo, so
  a forfeiture is an internal transfer rather than a loss to a counterparty. The
  incentive differs; the mechanism does not.
\end{enumerate}

Scheduling policy MAY restrict workloads by assurance class and by operator org
--- for instance, requiring \code{confidential} for a tenant's sensitive
inference, or restricting a regulated workload to a named org. That is a
constraint over fields, evaluated by the scheduler, and it is the one place the
difference is allowed to appear. It is not a second identity model.

\begin{remark}
No field marks a node as trusted for skipping a verification step, and none
should be added. Such a field is a target: whoever obtains it obtains the
assurance the network attaches to our own hardware, without the hardware. Hanzo
nodes reach their class through the same evidence as everyone else, which is
cheap for us because we control the hosts and can enable SecureBoot on them.
\end{remark}

% =====================================================================
\section{Residual risk}

What this design does not fix, stated plainly so nobody has to discover it later:

\begin{itemize}
\item \textbf{Class \code{self} is an honesty assumption.} Most of the fleet is
  at \code{self} today. For those nodes, every property in
  \S\ref{sec:notprove} stands, and the only controls are the bond and the
  challenge system. That is a real security posture --- economic rather than
  cryptographic --- and it must be described that way in anything a customer
  reads.
\item \textbf{Time-of-check to time-of-use.} A measurement describes a boot. A
  node measured honestly at $t$ can \code{kexec}, load a module, or write its
  root filesystem at $t + \epsilon$. Continuous attestation --- runtime
  measurement of what is loaded, not just what was --- is out of scope here and
  is the natural next specification.
\item \textbf{The bond bounds the loss, not the damage.} A node that leaks a
  tenant's data loses its bond; the tenant's data is still out. Confidentiality
  of tenant workloads on non-confidential VMs is a scheduling policy question,
  not an identity question.
\item \textbf{Sybil resistance is economic.} Nothing prevents one operator from
  running many nodes, each correctly identified. That is a feature; it becomes a
  problem only where a protocol counts nodes rather than weighing stake. No
  protocol specified here counts nodes.
\item \textbf{Vendor roots are trust anchors.} At class \code{platform} and above
  the verifier trusts a manufacturer's certificate chain and its revocation
  behavior. That is a smaller assumption than trusting every node operator, and
  it is not zero.
\item \textbf{The registration does not prove possession of the consensus key.}
  \emph{Observed}: the manager checks the BLS public key's length and the
  \code{nodeID}'s uniqueness, and verifies possession of neither. Uniqueness
  prevents two registrations of one \code{nodeID} on the same manager; it does
  not establish that the registrant holds the corresponding secrets. Requiring a
  proof of possession at \code{bind} time would close this; the primitive already
  exists (\code{luxfi/keys}, \code{pop.go}, \code{ProofOfPossession}).
\end{itemize}

% =====================================================================
\section{Implementation status and interface requests}
\label{sec:status}

\subsection{Implemented}

\begin{center}
\small
\begin{tabular}{@{}>{\raggedright\arraybackslash}p{0.28\textwidth}>{\raggedright\arraybackslash}p{0.64\textwidth}@{}}
\toprule
\textbf{Where} & \textbf{What} \\
\midrule
\code{hanzoai/node},\newline \code{hanzo-libs/hanzo-attest} &
\code{Boot}, canonical encoding, BLAKE3 \code{Measurement}; \code{Evidence} with
domain-separated \code{Kind}; \code{Key} (ed25519) as \code{Kind::Host} attester;
\code{Host} verifier; \code{Record} with \code{coherent()}; in-memory
\code{Registry}. \\[3pt]
\code{hanzoai/cli},\newline \code{src/commands/wallet.rs} &
Two custody models, no plaintext secrets: cloud custody where the key never
leaves KMS/MPC, and local custody in the OS keychain. EVM address derivation. \\[3pt]
\code{hanzoai/cloud},\newline \code{apps/wallet} &
\code{/v1/wallet/*} with four custody backends (\code{kms}, \code{mpc},
\code{treasury}, \code{safe}), key rotation, audit emission, a
\code{TierValidator} tier, and fail-closed behavior when the MPC cluster is
absent. \\[3pt]
\code{luxfi/kms} &
AES-256-GCM with ML-KEM-wrapped DEK, AAD bound to path/name/env/version, with a
test that a swapped DEK fails; hybrid X25519 + ML-KEM-768 transport. \\[3pt]
\code{luxfi/standard},\newline \code{contracts/validators} &
ACP-77 validator manager: bonded registration, \code{owner}, uptime-conditioned
rewards, churn limits, end-of-validation. \\[3pt]
\code{hanzoai/cloud},\newline \code{cli/link.go} &
\code{hanzo link}: starts \code{hanzod} via \code{hanzo node up} and registers
the host in the org's compute fleet, authorized by the IAM token. \\
\bottomrule
\end{tabular}
\end{center}

\subsection{Interface requests}
\label{sec:req}

These are changes in code owned elsewhere. They are requested here, with the
reason and the minimal shape, and are not made by this document.

\begin{description}
\item[IR-1 --- \code{hanzod}: the identity key is the wallet key, and is never
  plaintext.] \code{hanzo-bin/hanzo-node/src/utils/keys.rs} loads an ed25519
  secret from a plaintext \code{KEY=value} file or the environment;
  \code{src/runner.rs} writes it back with \code{std::fs::write}, at the umask.
  Requested: the node identity becomes a secp256k1 key whose address is $A_n$,
  obtained through the wallet custody path that already exists, so that
  \code{hanzod} holds an address and a signing capability rather than a secret.
  Rationale: \S\ref{sec:custody} C1--C2, and \S\ref{sec:gap} --- with a copyable
  key none of the properties in this document hold. If the change must land in
  two steps, the first step is the file mode and the second is the custody.

\item[IR-2 --- \code{hanzo-attest}: bind the preimage to chain, address and
  epoch.] \code{evidence::preimage} is
  $(\text{tag}, \text{kind}, \mu, \text{timestamp})$. Requested: add
  \code{chainID}, the signer's identity address, and the epoch; keep the existing
  kind byte. Rationale: without the address the statement is presentable under
  another registration; without the chain it crosses networks; without the epoch,
  equivocation (\S\ref{sec:equiv}) --- the fault that detects a copied key --- is
  undefined. Highest value of the five.

\item[IR-3 --- a node-binding registry contract.] The vendored manager preserves
  its ABI deliberately and should not be forked. Requested: the two-mapping,
  one-function contract of \S\ref{sec:stake}, so that \code{validationID} and
  $A_n$ agree on-chain. Optionally, require a proof of possession of $A_n$ at
  \code{bind}.

\item[IR-4 --- \code{hanzo-attest}: a TPM 2.0 kind, ahead of SNP and TDX.]
  \code{Kind} names \code{Snp} and \code{Tdx} and implements neither. On the
  measured fleet, two of three hosts have TPM 2.0 with a SHA-256 PCR bank and an
  event log, and zero have SNP or TDX. Requested: \code{Kind::Tpm} with a quote
  over PCRs and a verifier-supplied nonce, reporting $H(pk_n)$. It is the only
  one of the three kinds any host on the measured fleet can satisfy.

\item[IR-5 --- the epoch statement, separate from \code{Record}.]
  \code{Record} is per-machine
  and a node files many. Requested: one epoch statement per node per epoch, over
  $\mu_{\text{node}}$, doubling as the liveness signal so there is no second
  heartbeat protocol.
\end{description}

\subsection{Operational requests}

\begin{itemize}
\item Enable SecureBoot on the x86 host. It has a TPM 2.0 and SecureBoot is
  currently disabled, which caps it at class \code{self} even once IR-4 lands: an
  unenforced boot chain means the PCR values are chosen by whoever controls the
  boot.
\item Retire the local credential files under \code{\textasciitilde/.hanzo/} on
  every host that runs a node. Owner-only file permissions are L1 by the ladder
  of \S\ref{sec:ladder}; C1 requires L2 or better, and a node's credential is
  what C3 scopes.
\end{itemize}

% =====================================================================
\appendix

\section{The fleet, as measured}
\label{app:fleet}

Measured 4 September 2026, by direct inspection of the three hosts. Recorded
because the assurance classes in \S\ref{sec:class} are claims about real hardware
and should be checkable.

\begin{center}
\footnotesize
\begin{tabular}{@{}llll@{}}
\toprule
& \textbf{aarch64 host} & \textbf{x86 host} & \textbf{Apple-Silicon host} \\
\midrule
CPU & NVIDIA GB10 & AMD Ryzen AI MAX+ 395 & Apple M4 Max \\
Arch & aarch64 & x86\_64 & arm64 \\
OS & Ubuntu 24.04.4 & Ubuntu 24.04.4 & macOS 26.5.2 \\
Kernel & 6.17.0-1031-nvidia & 7.0.0-30-generic & Darwin \\
TPM & 2.0, SHA-256 PCR bank & 2.0 & none exposed \\
Event log & \code{binary\_bios\_measurements} & --- & --- \\
SecureBoot & enabled & \textbf{disabled} & n/a \\
SEV-SNP / TDX & n/a (aarch64) & \code{sev\_snp = N} & n/a \\
Monitor & KVM & cloud-hypervisor & Virtualization.framework \\
k3s & control-plane, v1.35.5+k3s1 & agent, v1.35.5+k3s1 & not joined \\
Class today & \code{self} (\code{platform} after IR-4) & \code{self} & \code{self} \\
\bottomrule
\end{tabular}
\end{center}

\noindent
The aarch64 host's PCR~0 read
\begin{center}\small\ttfamily
3d458cfe55cc03ea1f443f1562beec8df51c75e14a9fcf9a7234a13f198e7969
\end{center}
\noindent
which is what a class-\code{platform} quote would bind against.

The Apple-Silicon host carried \code{hanzo} 8.5.157 and \code{hanzo-vm} 2.0.0,
with OS assets under \code{\textasciitilde/.hanzo/vm} at version 2.0.0: kernel
\code{Image} (SHA-256 \code{2f71e062\dots{}3025}), \code{initramfs.cpio.gz}
(\code{397718d0\dots{}a6ff}), and a 4~GiB \code{rootfs.ext4}. Those are the files a
\code{Boot} measures --- Appendix~\ref{app:encoding} --- so the measurement is
reproducible by anyone holding the same release.

Two hosts are joined to one k3s cluster; the Mac is not currently a member. The
microVM is what would let it join on the same terms, which is the architectural
point of \S\ref{sec:class}: uniform interface, honestly non-uniform assurance.

\section{The boot encoding}
\label{app:encoding}

\emph{Observed}, \code{hanzo-attest/src/boot.rs}. Reproduced because
$\mu_{\text{node}}$ is what an epoch statement commits to, and a verifier must be
able to recompute it from the release.

\[
\mu = \mathrm{BLAKE3}\bigl(\mathrm{enc}(B)\bigr)
\]

where $\mathrm{enc}$ is: the domain tag \code{hanzo.boot.v1}, then the monitor
binary digest, the OS-asset version string, the kernel digest and the root
filesystem digest --- each length-prefixed with a big-endian \code{u64} --- then
\code{cpus}, \code{memory\_mb} and \code{disk\_mb} as big-endian \code{u64}, then
the argv count and each argument length-prefixed.

Three properties of the encoding, each with a test in the crate:

\begin{itemize}
\item \textbf{Paths do not enter.} Only digests. Moving an identical file does not
  change $\mu$.
\item \textbf{Length prefixes separate fields.} \code{["k3s","server"]} and
  \code{["k3sserver"]} measure differently.
\item \textbf{The kernel command line is pinned by the version.} It is built
  inside the monitor as a function of the OS-asset version, so the version stands
  for it without this crate restating a value it does not own. A verifier
  therefore checks the version against the published release; a monitor that
  changed the command line without changing the version would break the binding,
  and the OS-asset release process must not permit it.
\end{itemize}

\end{document}
